\chapter*{Résumé}
\addcontentsline{toc}{chapter}{Résumé}

Ce projet tuteuré a pour objectif d'analyser le parallèle entre le traitement du langage par les modèles de type Transformer et les mécanismes cognitifs du cerveau humain. Nous utilisons l'analyse de la dimensionnalité intrinsèque (ID) pour vérifier si l'IA effectue, comme le cerveau, une compression de l'information sémantique au cours de la lecture.

Pour répondre à cette problématique, nous comparons des phrases sémantiquement cohérentes à des phrases de type «~Jabberwocky~» (syntaxiquement correctes mais dénuées de sens). Nous mesurons l'évolution géométrique de ces représentations vectorielles couche par couche en implémentant l'algorithme TwoNN sur le modèle GPT-2. Les résultats obtenus confirment nos hypothèses : les phrases sémantiques créent une dynamique de «~ramping~» progressive dès les couches intermédiaires et un profil caractéristique de compression «~Hunchback~», tandis que le traitement du Jabberwocky présente une explosion de la dimensionnalité dans les couches supérieures, traduisant un échec d'unification.

Dans le premier chapitre, nous explorons l'état de l'art théorique reliant l'architecture des Transformers, le modèle neurolinguistique MUC et les propriétés géométriques de l'ID. Le deuxième chapitre détaille la méthodologie expérimentale et la génération automatique des stimuli. Le troisième chapitre présente la feuille de route initiale du projet. Enfin, le quatrième chapitre expose nos résultats empiriques, leurs analyses statistiques et leur discussion au regard des théories cognitives.